\begin{proof}
Fix a displayed trivializing open $U$ and let $x\in U$ be rigid.  By the
additional hypothesis, both $h_{x,\Omega_x}$ and its negative have positive
Hessian current.  Wanner's comparison theorem
\cite[Theorem~4.7 and Corollary~4.8]{Wanner19}, applied over the
algebraically closed complete field $\Omega_x$, makes both signs
Thuillier-subharmonic.  The proper connected curve
$D_{\Omega_x}^{\mathrm{an}}$ is compact and connected.  The maximum
principle \cite[Proposition~3.1.11]{Thu05}, applied to either sign at a
global extremum, shows that the corresponding extremum locus is open and
closed; hence $h_{x,\Omega_x}$ is constant.  The scalar-extension projection
\[
       D_{\Omega_x}^{\mathrm{an}}\longrightarrow
       D_{\mathscr H(x)}^{\mathrm{an}}
\]
is surjective by the surjectivity of analytic ground-field extension
\cite[\S1.3]{Berk12}; hence $h_x$ is constant.

We descend these fibrewise constants without assuming that $D$ has a
$K$-rational point.  Choose a closed point $p\in D$.  Its residue field
$L=\kappa(p)$ is a finite extension of $K$.  Put
$U_L=U\times_KL$ and let
\[
 q:U_L\longrightarrow U,
 \qquad
 s_p:U_L\longrightarrow U\times_KF
\]
be the finite faithfully flat projection and the multisection induced by
$p$.  In particular, $q$ is surjective.  Set
$b_L=h\circ s_p$.  On the strict space $U_L\times_UU_L$, the two pullbacks
of $b_L$ agree at every rigid point: its image in $U$ is rigid, and the two
values are values of the constant function on the corresponding analytic
fibre of $\pi$.  Rigid points are dense by
\cite[Proposition~2.1.15]{Berk12}, so the two pullbacks agree everywhere.
The finite surjection $q$ is closed and therefore a quotient map.  Thus
$b_L$ descends uniquely to a continuous function $b_U$ on $U$.

Pull the desired identity
\[
                         h=\operatorname{pr}_U^*b_U
 \label{eq:p062-scalar-stable-local-function-descent}
\]
back to $(U\times_KF)\times_UU_L$.  At every rigid point, both sides equal
the constant value on the rigid fibre over its image in $U$: the value of
$b_U$ is computed at any point of the nonempty finite fibre of $q$, where
it equals $b_L=h\circ s_p$.  Density of rigid points gives the pulled-back
identity everywhere, and surjectivity of the finite base change then gives
\eqref{eq:p062-scalar-stable-local-function-descent} on $U\times_KF$.

Choose a compactly supported ordinary CLD-smooth $(1,1)$-form $\eta$ on
$F$ with $\int_F\eta=1$.  Such a form is obtained from a nonzero
nonnegative compactly supported top form in a strict curve chart and
normalization; compare
\cite[Proposition~11.1 and Example~11.7]{GJR26}.  If $\beta$ is a compactly
supported ordinary smooth $(m-1,m-1)$-form on $U$, top fibre degree gives
\[
 \mathrm d'\mathrm d''\bigl(\operatorname{pr}_U^*\beta
       \wedge\operatorname{pr}_F^*\eta\bigr)
 =\operatorname{pr}_U^*(\mathrm d'\mathrm d''\beta)
       \wedge\operatorname{pr}_F^*\eta.
\]
Pairing with the zero current of $h$, using
\eqref{eq:p062-scalar-stable-local-function-descent}, and integrating over
the fibre yields
\[
 0=\int_{U\times_KF}h\,
       \operatorname{pr}_U^*(\mathrm d'\mathrm d''\beta)
       \wedge\operatorname{pr}_F^*\eta
   =\int_U b_U\,\mathrm d'\mathrm d''\beta.
\]
Here integration is over the fixed $K$-analytic factor $F$, so the
normalization is literally $\int_F\eta=1$; no scalar-extension
normalization is being used.  For completeness, choose a smooth cutoff
equal to one near the support of the displayed test form, approximate the
resulting compactly supported continuous multiplier uniformly by ordinary
smooth functions, apply the smooth projection formula
\cite[Theorem~16.5.4 and Corollary~16.5.5]{CLD25}, and pass to the limit
against the fixed compactly supported signed measure.  This proves the
continuous-multiplier fibre-integration identity used above.  It follows that
$\mathrm d'\mathrm d''[b_U]=0$.  The identical calculation with a
compactly supported ordinary smooth $(m-1,m)$-form proves
$\mathrm d'[b_U]=0$ when $\mathrm d'[h]=0$.

On overlaps, the local pullback identities and the nonempty fibres force
the functions $b_U$ to agree.  They glue to a unique continuous function
$b$ on $B$.  Since current differentiation and restriction to an ordinary
open subspace commute by their test-form definitions, the local current
identities glue, proving
\eqref{eq:p062-scalar-stable-curve-descent} and the first-current assertion.
\end{proof}
